\chapter[PITR, WAL y Backup]{\emj{💾}~PITR, WAL y Backup}

\section[Write-Ahead Log en profundidad]{\emj{📝}~Write-Ahead Log en profundidad}

WAL (Write-Ahead Log) es el mecanismo de durabilidad de PostgreSQL.
\textbf{Regla fundamental:} ningún cambio se escribe al heap (datos) sin antes haberse escrito al WAL.
Esto garantiza que tras un crash, el WAL puede re-aplicarse para recuperar el estado consistente.

\begin{teoriabox}{Arquitectura del WAL}
\begin{itemize}
  \item \textbf{WAL segment:} archivos de 16MB en \texttt{pg\_wal/}. Nombrados por LSN (Log Sequence Number): \texttt{000000010000000000000001}.
  \item \textbf{LSN (Log Sequence Number):} posición única en el flujo WAL. Formato: \texttt{X/YYZZZZZZ}. Siempre creciente.
  \item \textbf{Checkpoint:} evento periódico que garantiza que todas las páginas sucias hasta cierto LSN están escritas al heap. Permite reutilizar WAL anterior al checkpoint.
  \item \textbf{WAL writer:} proceso que flushea WAL buffers al disco. La frecuencia afecta latencia de commits.
  \item \textbf{WAL sender:} envía WAL a standbys (replicación) o al archivador.
\end{itemize}
\end{teoriabox}

\begin{lstlisting}[caption={Parámetros WAL críticos}]
-- postgresql.conf:
wal_level = replica          -- mínimo para replicación y archivado
                             -- logical: replicación lógica también
checkpoint_timeout = 5min    -- checkpoint automático cada X (default 5min)
checkpoint_completion_target = 0.9  -- escribir páginas sucias en 90% del intervalo
max_wal_size = 1GB           -- WAL puede crecer hasta X entre checkpoints
min_wal_size = 80MB          -- retener al menos X de WAL
wal_compression = on         -- comprimir WAL (reduce I/O a costo de CPU mínimo)

-- Ver LSN actual y estado de WAL:
SELECT pg_current_wal_lsn();               -- posición actual
SELECT pg_walfile_name(pg_current_wal_lsn());  -- nombre de archivo actual
SELECT pg_wal_lsn_diff(pg_current_wal_lsn(), '0/0');  -- bytes totales

-- Ver checkpoints recientes:
SELECT * FROM pg_stat_bgwriter;
-- checkpoints_timed: cuántos por timeout (bueno)
-- checkpoints_req: cuántos forzados por WAL lleno (malo → aumentar max_wal_size)
\end{lstlisting}

\section[WAL Archiving]{\emj{📦}~WAL Archiving}

El archivado de WAL copia cada segmento completado a un almacenamiento externo (S3, NFS, etc.).
Combinado con un base backup, permite PITR a cualquier punto en el tiempo.

\begin{lstlisting}[caption={Configurar WAL archiving}]
-- postgresql.conf:
archive_mode = on
archive_command = 'aws s3 cp %p s3://mi-bucket/wal-archive/%f'
-- %p = ruta completa del archivo WAL
-- %f = nombre del archivo (sin ruta)
-- El comando debe retornar 0 solo si la copia fue exitosa

-- Alternativa con rsync:
-- archive_command = 'rsync -a %p backup-server:/mnt/wal-archive/%f'

-- archive_cleanup_command: limpiar WAL archivados ya no necesarios (con pgBackRest):
-- archive_cleanup_command = 'pg_archivecleanup /mnt/wal-archive %r'

-- Verificar que el archivado funciona:
SELECT archived_count, last_archived_wal, last_archived_time,
       failed_count, last_failed_wal, last_failed_time
FROM pg_stat_archiver;
-- failed_count > 0 y creciendo: problema con archive_command
\end{lstlisting}

\section[pgBackRest: backup empresarial]{\emj{🛡️}~pgBackRest: backup empresarial}

pgBackRest es la herramienta de backup más completa para PostgreSQL: maneja base backups, WAL archiving, compresión, cifrado y PITR.

\begin{lstlisting}[language=, caption={pgbackrest.conf — configuración básica}]
[global]
repo1-path=/var/lib/pgbackrest
repo1-retention-full=2          ; retener 2 full backups
repo1-retention-diff=7          ; retener 7 differential backups
repo1-cipher-type=aes-256-cbc
repo1-cipher-pass=supersecret

# Para S3:
# repo1-type=s3
# repo1-s3-bucket=mi-bucket-backups
# repo1-s3-region=us-east-1

[global:archive-push]
compress-level=3

[mydb]
pg1-path=/var/lib/postgresql/data
pg1-port=5432
pg1-user=postgres
\end{lstlisting>

\begin{lstlisting}[language=, caption={pgBackRest — operaciones comunes}]
# Inicializar el repositorio:
pgbackrest --stanza=mydb stanza-create

# Full backup:
pgbackrest --stanza=mydb --type=full backup

# Differential backup (desde último full):
pgbackrest --stanza=mydb --type=diff backup

# Incremental backup (desde último backup cualquiera):
pgbackrest --stanza=mydb --type=incr backup

# Ver información de backups:
pgbackrest --stanza=mydb info

# Restaurar último backup:
pgbackrest --stanza=mydb restore

# PITR: restaurar a punto específico en el tiempo:
pgbackrest --stanza=mydb restore \
  --type=time "--target=2025-06-15 14:30:00"

# Verificar integridad del repositorio:
pgbackrest --stanza=mydb check
\end{lstlisting}

\section[Procedimiento PITR]{\emj{🕰️}~Procedimiento PITR}

Point-In-Time Recovery: restaurar la base a cualquier momento entre el último base backup y el WAL archivado más reciente.

\begin{lstlisting}[caption={PITR manual con pg\_restore + WAL}]
-- ESCENARIO: DROP TABLE accidental a las 14:23:00
-- Base backup más reciente: 02:00 AM
-- WAL archivado hasta: 14:30 (5 min después del desastre)

-- PASO 1: Detener PostgreSQL
-- systemctl stop postgresql

-- PASO 2: Mover datos actuales:
-- mv /var/lib/postgresql/data /var/lib/postgresql/data.broken

-- PASO 3: Restaurar base backup:
-- pgbackrest --stanza=mydb restore --delta  (o desde tar manual)

-- PASO 4: Configurar recovery target en postgresql.conf:
-- (PG 12+: en el mismo postgresql.conf)
restore_command = 'pgbackrest --stanza=mydb archive-get %f "%p"'
recovery_target_time = '2025-06-15 14:22:59'  -- un segundo antes del DROP
recovery_target_action = promote  -- promote: cuando llegue al target, se convierte en primary

-- PASO 5: Crear señal de recovery:
-- touch /var/lib/postgresql/data/recovery.signal

-- PASO 6: Iniciar PostgreSQL:
-- systemctl start postgresql
-- PG aplica WAL hasta 14:22:59 y promueve

-- PASO 7: Verificar que la tabla existe y datos están correctos
-- PASO 8: Tomar nuevo base backup
\end{lstlisting}

\begin{lstlisting}[caption={RPO y RTO — métricas de recuperabilidad}]
-- RPO (Recovery Point Objective): ¿cuántos datos puedes perder?
-- Con archivado WAL cada ~60s: RPO ~= 60 segundos de datos
-- Con replicación síncrona: RPO = 0

-- RTO (Recovery Time Objective): ¿cuánto tiempo toma recuperarse?
-- Depende de: tamaño del backup + cantidad de WAL a replay
-- pgBackRest con backup diferencial puede reducir RTO vs full backup

-- Calcular WAL generado por día (útil para sizing de S3):
SELECT pg_size_pretty(SUM(size))
FROM pg_ls_waldir();

-- Probar recovery ANTES de necesitarla (no en producción):
-- Montar un servidor de prueba, restaurar, validar datos críticos
-- El backup que nunca probaste no es un backup
\end{lstlisting}

\begin{entrevistabox}
\begin{enumerate}
  \item ¿Qué es un LSN y cómo se relaciona con el WAL?
  \item Explica la diferencia entre un checkpoint y un backup.
  \item ¿Qué es PITR? ¿Qué necesitas para poder hacerlo?
  \item Si \texttt{checkpoints\_req} está creciendo en \texttt{pg\_stat\_bgwriter}, ¿qué indica y cómo lo corriges?
  \item ¿Cuál es la diferencia entre RPO y RTO? Da un ejemplo concreto.
  \item Un DROP TABLE ocurrió en producción hace 10 minutos. Tienes base backup y WAL archivado. Describe los pasos de recuperación.
  \item ¿Por qué ``el backup que nunca probaste no es un backup''?
\end{enumerate}
\end{entrevistabox}
